\section{Conclusion}
This paper studies the portfolio consequences of replacing heterogeneous signal uncertainty with a common scalar adjustment. The industry replication recovers the performance improvement associated with enhanced portfolio optimization and documents its redistribution of risk away from small-eigenvalue directions. In the known-uncertainty experiment, the best scalar rule loses 12.4\% and 28.7\% of mean heterogeneous utility for the two nonuniform profiles. The loss grows with exposure alignment and remains after the scalar parameter has been optimized.

The feasible extension gives a different empirical answer. A bootstrap-based relative uncertainty profile yields annual net certainty equivalent differences of 27.65 basis points in 1970--2018 and $-98.52$ in 2019--2025, with both confidence intervals including zero. Higher mean returns are accompanied by additional variance and costs. These results establish the value of uncertainty differentiation in the specified known-input model, but do not establish an improvement in historical net utility from the implemented estimator.

The scope of this evidence is limited by the experimental and empirical designs. The simulation uses a Gaussian model, one fixed uncertainty profile per dispersion level and known covariance; its Monte Carlo intervals condition on those profiles. The calibration experiment uses a correctly specified linear conditional mean and does not reproduce the full construction of historical momentum. The historical extension uses one signal, one industry dataset and a diagonal uncertainty adjustment, with common rather than separately optimized tuning. The 2019--2025 period contains only 84 monthly observations. Return-block intervals condition on estimated forecasts and selection histories, so they do not capture uncertainty from rebuilding the entire research procedure. Industry portfolio returns and constant trading-cost assumptions also abstract from security-level execution. These features define the population and implementation to which the reported comparisons apply.

The main finding is therefore conditional and economic: uncertainty differences can matter for optimal allocation even when one scalar parameter is chosen well, but realizing that value requires sufficiently informative estimates and a portfolio policy that preserves the gain after risk and trading costs. The replication, simulation and matched evaluation together identify that distinction within a single, reproducible framework.
